\chapter{环境要求与安装部署}
\label{ch:2}

本章内容：说明 EdgeFlow v0.7.0 的运行环境要求，并给出完整的安装部署步骤，包括获取软件、使用 \texttt{keadm} 初始化云端与接入边缘节点、Helm 部署，以及部署验证、卸载清理与升级回滚。

\section{环境要求}

\begin{tabularx}{\textwidth}{|l|X|}
  \hline
  \textbf{部署角色} & \textbf{环境要求} \\
  \hline
  云端（cloudcore） & 支持 \texttt{linux/amd64} 与 \texttt{linux/arm64}；建议 2 核 4 GB 以上；建议配置持久化存储：embed 模式使用数据目录（Helm 默认 PVC 1Gi，emptyDir 会丢数据），外部 etcd 模式需共享 etcd 集群 \\
  \hline
  Kubernetes 集群 & 使用 \texttt{keadm} 产物或 Helm 部署时需准备集群；需要 \texttt{kubectl}（与集群版本兼容）与 Helm v3+ \\
  \hline
  边缘节点（edgecore） & Linux 系统（支持 systemd）；Docker daemon 运行中（Edged 的容器运行时）；可访问云端 CloudHub 端口（默认 10000，NodePort 部署时为集群分配的节点端口）；SQLite 数据目录可写 \\
  \hline
  管理机（keadm） & 任意 amd64/arm64 主机，仅用于生成离线部署产物，不直接操作集群 \\
  \hline
  源码构建 & Go 1.26+ 与 Make；构建容器镜像还需 Docker \\
  \hline
  可选依赖 & mosquitto（MQTT 数据面）；缺失时主链路不受影响 \\
  \hline
  外部 etcd 集群（可选） & 使用外部 etcd 模式（v0.5.0 起）时需准备：3 节点奇数、同地域、quota 256MiB、compaction 1h；生产建议启用 TLS/mTLS（非回环端点且无 TLS 时云端拒绝启动） \\
  \hline
\end{tabularx}

\section{获取软件}

\subsection{方式一：源码构建}

在仓库根目录执行：

\begin{lstlisting}
make build
\end{lstlisting}

构建产物位于 \texttt{bin/} 目录：\texttt{bin/cloudcore}、\texttt{bin/edgecore}。\texttt{keadm} 可用 \texttt{go build -o bin/keadm ./cmd/keadm} 构建；\texttt{mock-cloudhub} 为联调工具，用 \texttt{go run ./hack/mock-cloudhub} 运行（无需构建）。

\subsection{方式二：使用发布制品}

发布目录 \texttt{release/v0.7.0/} 提供预编译二进制：\texttt{cloudcore}、\texttt{edgecore}、\texttt{keadm} 每组件均含 \texttt{linux-amd64}、\texttt{linux-arm64} 与 \texttt{darwin-arm64} 三个平台版本（共 9 个二进制），以及 Helm Chart 包 \texttt{edgeflow-0.7.0.tgz}（历史版本制品见 \texttt{release/} 对应子目录）。下载后可使用 \texttt{checksums.txt} 校验文件完整性，并将二进制放入 \texttt{PATH} 或仓库 \texttt{bin/} 目录后按需执行。

\subsection{方式三：容器镜像}

可自行构建镜像，或直接使用已发布的 \texttt{edgeflow/cloudcore:v0.7.0}、\texttt{edgeflow/edgecore:v0.7.0}。本地构建命令（仓库根目录执行）：

\begin{lstlisting}
docker build -f build/docker/Dockerfile --target cloudcore -t edgeflow/cloudcore:v0.7.0 .
docker build -f build/docker/Dockerfile --target edgecore -t edgeflow/edgecore:v0.7.0 .
\end{lstlisting}

\section{安装云端（cloudcore）}

\subsection{方式一：keadm 生成 Kubernetes 部署产物（推荐）}

\texttt{keadm init} 在管理机上生成云端部署产物，产物可提交到 Kubernetes 集群执行。参数如下：

\begin{tabularx}{\textwidth}{|l|l|X|}
  \hline
  \textbf{参数} & \textbf{默认值} & \textbf{说明} \\
  \hline
  \texttt{--cloudcore-image} & \texttt{edgeflow/cloudcore:v0.7.0} & cloudcore 容器镜像 \\
  \hline
  \texttt{--tls} & 关 & 启用云边 mTLS，注入 TLS 相关环境变量 \\
  \hline
  \texttt{--tls-san} & 空 & 证书 SAN（Subject Alternative Name），逗号分隔，如 \texttt{IP:1.2.3.4,DNS:host}；仅 \texttt{--tls} 时生效 \\
  \hline
  \texttt{--service-type} & \texttt{NodePort} & Service 类型：\texttt{NodePort}（边缘跨集群接入）或 \texttt{ClusterIP}（仅集群内访问） \\
  \hline
  \texttt{--output-dir} & \texttt{./keadm-out} & 产物输出目录 \\
  \hline
\end{tabularx}

生成命令示例：

\begin{lstlisting}
# 最简（明文通道，本地/测试用）
keadm init --output-dir=./keadm-out

# 生产推荐：启用云边 mTLS，并注入证书 SAN（边缘节点用 IP 接入时必须覆盖访问地址）
keadm init --tls --tls-san=IP:192.168.1.10 --output-dir=./keadm-out
\end{lstlisting}

产物说明：

\begin{itemize}
  \item \texttt{cloudcore.yaml}：包含 Deployment（副本 1、探针指向 \texttt{/healthz}、安全上下文 nonroot）与 Service（\texttt{NodePort} 类型，暴露 HTTP 8080 与 CloudHub 10000 两个端口；hub 端口由集群自动分配在 30000--32767）；
  \item \texttt{NOTES.txt}：部署步骤、验证方法、Helm 替代路径与 mTLS 说明。
\end{itemize}

在集群上执行产物：

\begin{lstlisting}
kubectl apply -f cloudcore.yaml

# 验证
kubectl get deploy,svc,pods -l app.kubernetes.io/component=cloudcore
kubectl port-forward svc/edgeflow-cloudcore 8080:8080
curl http://127.0.0.1:8080/healthz        # 期望 HTTP 200

# 获取边缘节点接入用的 CloudHub 节点端口
kubectl get svc edgeflow-cloudcore -o jsonpath='{.spec.ports[?(@.name=="hub")].nodePort}'
\end{lstlisting}

\subsection{方式二：Helm 部署}

使用仓库内置 Chart 部署云端：

\begin{lstlisting}
helm install edgeflow build/charts/edgeflow

# 集群外边缘节点接入时：开启 hub 端口并改用 NodePort
helm install edgeflow build/charts/edgeflow \
  --set service.hubEnabled=true --set service.type=NodePort
\end{lstlisting}

验证部署：

\begin{lstlisting}
kubectl get deploy,svc,pods -l app.kubernetes.io/instance=edgeflow
kubectl port-forward svc/edgeflow-cloudcore 8080:8080
curl http://127.0.0.1:8080/healthz
\end{lstlisting}

关键配置项：\texttt{cloudcore.image.repository/tag}（镜像地址与版本，默认 \texttt{edgeflow/cloudcore:v0.7.0}）、\texttt{service.type}（默认 \texttt{ClusterIP}）、\texttt{service.hubEnabled}（默认 \texttt{false}，集群外边缘接入需开启并配合 NodePort/LoadBalancer）。

\textbf{v0.4.0 起新增关键配置项}（云端持久化与高可用形态）：

\begin{itemize}
  \item \texttt{cloudcore.replicaCount}：默认 1。embed 模式必须为 1——多副本各自内嵌 etcd 会脑裂，Chart 渲染期 \texttt{\{\{ fail \}\}} 守卫直接报错；外部 etcd 模式（v0.6.0 起）可 \texttt{>1}，此时自动注入 \texttt{EDGEFLOW\_CLOUDCORE\_MULTI\_REPLICA=1}（真多活）；
  \item \texttt{cloudcore.etcd.enabled}：默认 true（注入 \texttt{EDGEFLOW\_CLOUDCORE\_ETCD\_ENABLED=true}）。设为 false 时强制纯内存模式并忽略 \texttt{external.*}（数据重启丢失，仅限测试）；
  \item \texttt{etcd.persistence.enabled/size}：默认 true/1Gi，PVC 挂载 \texttt{/data}。改为 emptyDir 后 etcd 数据不落盘，持久化名存实亡，生产环境请保持默认；
  \item \texttt{etcd.external.enabled/endpoints}：v0.5.0 起外部 etcd 模式——enabled 为 true 且 endpoints 非空时注入 \texttt{EDGEFLOW\_CLOUDCORE\_ETCD\_ENDPOINTS}（逗号连接），跳过 embed、不创建 PVC；endpoints 为空时渲染失败；
  \item \texttt{etcd.external.tls.\{ca,cert,key\}} 与 \texttt{etcd.external.allowInsecure}：v0.5.0 起——非回环端点且未配置 TLS 时云端拒绝启动（明文护栏），仅限可信内网可用 \texttt{allowInsecure=true} 逃生；
  \item \texttt{etcd.external.nodeLeaseTTL}：v0.6.0——非空时注入 \texttt{EDGEFLOW\_CLOUDCORE\_NODE\_LEASE\_TTL}（默认 300s，外部模式判活阈值）；
  \item \texttt{cloudcore.resources}：建议 requests 256Mi / limits 1Gi（v0.4.0 起二进制体积与常驻内存上升）。
\end{itemize}

\subsection{方式三：裸进程运行（单机联调）}

开发或验证环境可直接运行二进制：

\begin{lstlisting}
./bin/cloudcore
\end{lstlisting}

默认监听 REST API 8080 与 CloudHub 10000 端口；端口可通过配置文件（\texttt{config/cloudcore.json}，可用 \texttt{--config} 指定路径）、环境变量或命令行覆盖，优先级为命令行 > 环境变量 > 配置文件 > 默认值（详见第 3 章 \autoref{ch:3}）。此方式仅用于单机联调，生产环境请使用方式一或方式二。

\section{接入边缘节点（edgecore）}

\texttt{keadm join} 在管理机或边缘节点上生成边缘接入产物。参数如下：

\begin{tabularx}{\textwidth}{|l|l|X|}
  \hline
  \textbf{参数} & \textbf{默认值} & \textbf{说明} \\
  \hline
  \texttt{--cloudcore-ip} & 必填 & 云端 CloudHub 节点 IP（IPv4/IPv6 均可） \\
  \hline
  \texttt{--cloudcore-port} & \texttt{10000} & CloudHub 端口；NodePort 部署时填集群分配的节点端口 \\
  \hline
  \texttt{--token} & 必填 & 接入令牌：写入 \texttt{edgecore.env} 并由 edgecore 注册时携带；云端设 \texttt{EDGEFLOW\_CLOUDCORE\_NODE\_TOKEN} 后启用校验 \\
  \hline
  \texttt{--node-id} & \texttt{edge-<主机名>} & 边缘节点 ID，仅允许字母、数字、点、连字符与下划线（\texttt{\textasciicircum{}[A-Za-z0-9.\_-]+\$}） \\
  \hline
  \texttt{--tls} & 关 & 启用 mTLS，地址自动使用 \texttt{wss://} \\
  \hline
  \texttt{--output-dir} & \texttt{./keadm-out} & 产物输出目录 \\
  \hline
\end{tabularx}

生成命令示例：

\begin{lstlisting}
keadm join --cloudcore-ip=192.168.1.10 --token=<token> --output-dir=./keadm-out

# TLS 集群 + NodePort 部署（端口为集群分配的 hub 节点端口）
keadm join --cloudcore-ip=192.168.1.10 --cloudcore-port=31000 \
  --token=<token> --node-id=edge-worker-01 --tls --output-dir=./keadm-out
\end{lstlisting}

产物说明：

\begin{itemize}
  \item \texttt{edgecore.env}：环境变量文件，键名与 edgecore 读取的环境变量一一对应，包括 \texttt{EDGEFLOW\_EDGECORE\_NODE\_ID}、\texttt{EDGEFLOW\_EDGECORE\_CLOUD\_ADDR}（\texttt{ws://<ip>:<port>/v1/edge}）、\texttt{EDGEFLOW\_EDGECORE\_DB\_PATH}（\texttt{/var/lib/edgeflow/edgeflow.db}）、\texttt{EDGEFLOW\_EDGECORE\_MQTT\_ADDR}（\texttt{tcp://127.0.0.1:1883}）等；\texttt{--tls} 时追加 TLS 与证书目录配置；\texttt{EDGEFLOW\_EDGECORE\_TOKEN} 为接入令牌键（edgecore 注册时携带，云端设 \texttt{EDGEFLOW\_CLOUDCORE\_NODE\_TOKEN} 后启用校验）；
  \item \texttt{edgecore.service}：systemd 单元文件，通过 \texttt{EnvironmentFile=/etc/edgeflow/edgecore.env} 加载配置，\texttt{ExecStart=/usr/local/bin/edgecore}，崩溃自动重启；
  \item \texttt{install.sh}：一键安装脚本（需 root），安装二进制、环境文件与 systemd 单元并启用服务；
  \item \texttt{README.md}：接入说明与手动安装片段。
\end{itemize}

在边缘节点上执行：

\begin{lstlisting}
# 1. 将产物与 edgecore 二进制（与节点架构匹配）拷贝到边缘节点
# 2. 执行安装脚本（需 root）
sudo ./install.sh

# 3. 验证
systemctl status edgecore
journalctl -u edgecore -f
\end{lstlisting}

mTLS 说明：启用 \texttt{--tls} 后，edgecore 首次运行会在 \texttt{/etc/edgeflow/certs} 自动生成或加载证书（幂等）；云端侧需保证服务端证书 SAN 覆盖边缘节点访问的地址，否则云边连接会被拒绝。

\section{部署验证}

部署完成后，建议按以下清单验证：

\begin{itemize}
  \item 云端 \texttt{/healthz} 返回 HTTP 200 与版本信息；
  \item 边缘节点已注册：云端查询 \texttt{/api/v1/nodes} 能看到节点且状态为 \texttt{Ready}（查询方法见第 4 章 \autoref{ch:4}）；
  \item \texttt{systemctl status edgecore} 显示 active (running)，日志无致命错误；
  \item 可下发一个测试 Pod 验证端到端链路（见第 3 章 \autoref{ch:3}）；
  \item 配置热重载：cloudcore 与 edgecore 支持 SIGHUP 触发与每 60 秒轮询两种热重载方式（云端 HTTP 端口可热切换，部分配置项变更需重启，生效边界见第 5 章 \autoref{ch:5}）。
\end{itemize}

\section{卸载与清理}

\begin{lstlisting}
# 清理 keadm 生成的产物（交互确认；--force 跳过确认，幂等）
keadm reset --output-dir=./keadm-out

# 卸载 Helm 部署
helm uninstall edgeflow

# 清理 Edged 管理的容器（按标签精确删除，不影响其他容器）
docker rm -f $(docker ps -aq --filter label=edgeflow.pod) 2>/dev/null

# 清理本地 SQLite 元数据（生产环境请按备份策略处理）
rm -f data/edgeflow.db
\end{lstlisting}

\section{升级与回滚}

\texttt{keadm} 在产物层面提供升级与回滚能力（执行前自动备份并记录操作台账）：

\begin{lstlisting}
# 升级产物到新版本（操作人经环境变量传入，写入操作台账）
export KEADM_OPERATOR=alice
keadm upgrade --version=v0.7.0 --output-dir=./keadm-out

# 回滚到最近一次备份
keadm rollback --latest --output-dir=./keadm-out

# 查询操作台账
keadm ops-ledger --limit=10
\end{lstlisting}

升级或回滚后需重新应用产物（\texttt{kubectl apply} 或 \texttt{./install.sh}）。注意：v0.4.0 起云端为分级持久化——注册台账与设备 Desired 跨重启保留，重启后节点短暂 Unknown，边缘节点下一次心跳（默认 30s）即翻新为 Ready；Pod 状态与上报数据重启后 $\le$1 个上报周期自愈。云端多副本（外部 etcd 模式）升级/回滚务必全停再全起，禁止混合版本混跑（详见第 7 章 \autoref{ch:7}）。

\section{常见问题}

\begin{tabularx}{\textwidth}{|l|X|}
  \hline
  \textbf{现象} & \textbf{处理} \\
  \hline
  \texttt{缺少必填参数 --cloudcore-ip} 或 \texttt{--token} & 补全对应参数后重试 \\
  \hline
  \texttt{--node-id 含空白字符} & 使用 \texttt{--node-id=edge-xxx} 显式指定合法 ID \\
  \hline
  \texttt{kubectl apply} 报 schema 校验失败 & 检查 kubectl 版本与集群 API 版本是否兼容 \\
  \hline
  edgecore 起不来 & \texttt{journalctl -u edgecore -e} 查看日志；确认可达 \texttt{--cloudcore-ip} 的 hub 端口；确认 env 文件键值未被手改 \\
  \hline
  云边连接被拒（TLS） & 证书 SAN 未覆盖访问地址，云端以 \texttt{--tls-san=IP:<访问IP>} 重新 init 并 apply \\
  \hline
  Pod 容器未创建 & 确认 Docker daemon 运行（\texttt{docker info}）；镜像本地未缓存时首次拉取需要时间 \\
  \hline
  云端重启后节点仍可见但状态 Unknown & v0.4.0 起注册台账跨重启保留，属正常现象：等待边缘节点下一次心跳（默认 30s）后自动翻新为 Ready，无需重新注册；Pod 状态与上报数据同样 $\le$1 个上报周期自愈 \\
  \hline
  embed 模式部署多副本失败（渲染报错） & Chart 内置 \texttt{\{\{ fail \}\}} 守卫：embed 模式 \texttt{replicaCount} 必须为 1（多副本各自内嵌 etcd 会脑裂）；需要多副本请改用外部 etcd 模式（v0.6.0 起支持） \\
  \hline
  外部 etcd 模式启动失败 & 按启动日志区分三类原因：探活失败（Unavailable，地址/连通性错误）、鉴权被拒（PermissionDenied，证书或 etcd 权限角色问题）、明文护栏（非回环端点且未配置 TLS 被拒绝启动，需配置 TLS 或仅在可信内网用 \texttt{ETCD\_ALLOW\_INSECURE=1} 逃生）；详细排障见部署文档 §10.7.6 \\
  \hline
\end{tabularx}
